\abstract{%
Documents (paper media, images, or electronic files containing textual and graphical information) are ubiquitous in daily office work, online dissemination, and governmental and enterprise workflows; automatically parsing, retrieving, and supporting decision-making based on their content constitutes a core requirement of Document Intelligence. In real-world settings, many documents are first converted into images via scanning or photographing; consequently, document processing typically starts from visual inputs: on the one hand, the system must accurately localize layout elements such as text blocks, tables, figures, and headings, and perform text detection and recognition; on the other hand, it must conduct higher-level semantic understanding and reasoning based on structured inputs to answer queries, extract key information, and generate verifiable results. With advances in deep learning and large language models, the research focus has gradually expanded from ``character-level recognition'' to ``document-level understanding and question answering'' across regions, pages, and modalities, and the field has correspondingly developed evaluation benchmarks that emphasize end-to-end capability and robustness.

Following this evolution, this paper organizes the survey in the order of ``perception first, then reasoning, and finally end-to-end unified modeling,'' and aligns evaluation suites with methodological lineages. Part~I (Document Layout Analysis and OCR) focuses on layout parsing starting from page geometry, including detection/segmentation-based layout element recognition, layout-aware models that incorporate layout information into pretrained representation learning, and robustness and cross-domain generalization under real-world document distributions; it further reviews key technical points of OCR from traditional pipelines to deep learning and lightweight deployment, emphasizing error propagation and engineering constraints when OCR serves as the ``entry point'' for downstream understanding tasks. Part~II (Document Understanding and Question Answering) systematically summarizes three mainstream scenarios built on structured inputs: structure-aware representations for tables, modular/executable reasoning and verifiable question answering; retrieval-augmented generation (RAG) and reasoning-driven retrieval for long texts; and multimodal pretraining and instruction alignment for visually rich documents, along with the evolution toward OCR-free modeling and multi-page, long-context understanding. Finally, we summarize the datasets and benchmarks associated with these two stages, covering task settings from single-page to multi-page documents and from closed sets to real-world enterprise documents, providing reproducible baselines for comparison and a systematic research roadmap for future studies.
}
